\section{A distinct noncanonical connection API}
\label{sec:connection}

\subsection{The minimal compiled API}

The compiled Chapter III API assumes only normalization at each chart.  All
chart transports act on one module $M$, so chart labels distinguish routes
but not Lean types.

\begin{definition}[Noncanonical connection normalized at each chart]
\label{def:minimal-connection}
A minimal connection is a family
\[
\Gamma_{b\leftarrow a}\in\operatorname{Aut}_R(M)
\qquad(a,b\in B)
\]
with only
\[
\Gamma_{a\leftarrow a}=\id_M.
\]
Neither inversion on reversed edges nor the cocycle law is a field of the
structure.
\end{definition}
\statusleanthree{Komorebi.PaperIII.Transport.Connection}

At a triple $(a,b,c)$, the compiled cocycle predicate is
\[
\Gamma_{c\leftarrow b}\Gamma_{b\leftarrow a}
=\Gamma_{c\leftarrow a}.
\]
Flatness means this equality holds for every triple.  The formal
multiplicative relative defect and triangle holonomy are
\[
D^\times_{c,b,a}
=
\bigl(\Gamma_{c\leftarrow b}\Gamma_{b\leftarrow a}\bigr)
\Gamma_{c\leftarrow a}^{-1},
\qquad
H_{c,b,a}
=
\Gamma_{a\leftarrow c}
\Gamma_{c\leftarrow b}
\Gamma_{b\leftarrow a}.
\]
The defect is based at the target $c$.  The holonomy is based at $a$.
\statusleanthree{Komorebi.PaperIII.Transport.Connection.CocycleAt; Komorebi.PaperIII.Transport.Connection.Flat; Komorebi.PaperIII.Transport.Connection.cocycleDefect; Komorebi.PaperIII.Transport.Connection.triangleHolonomy}

\begin{theorem}[Checked flatness consequences]
\label{thm:minimal-flatness}
For the minimal connection API, flatness implies
\[
D^\times_{c,b,a}=\id_M,
\qquad
H_{c,b,a}=\id_M
\]
for every $a,b,c$.
\end{theorem}
\statusleanthree{Komorebi.PaperIII.Transport.Connection.cocycleDefect_eq_refl_of_flat; Komorebi.PaperIII.Transport.Connection.triangleHolonomy_eq_refl_of_flat}

\begin{proof}
Flatness replaces the two-edge route in $D^\times$ by the direct route,
which cancels its inverse.  For $H$, flatness first collapses $a\to b\to c$
to $a\to c$, then collapses $a\to c\to a$ to the normalized transport
$a\to a$.
\end{proof}

\begin{proposition}[The minimal API admits nontrivial holonomy]
\label{prop:twisted-loop}
There is a connection on three points over $\mathbb Q$ for which transport
$0\to1$ is multiplication by $2$, every other transport between distinct
endpoints is the identity, and
\[
H_{2,1,0}(1)=2.
\]
In particular $H_{2,1,0}\ne\id$.
\end{proposition}
\statusleanthree{Komorebi.PaperIII.Transport.twistedThreePointConnection; Komorebi.PaperIII.Transport.twisted_triangleHolonomy_apply_one; Komorebi.PaperIII.Transport.twisted_triangleHolonomy_ne_refl}

\begin{boundary}[Why the minimal defect and loop are distinct]
Without reverse-edge inversion,
$\Gamma_{a\leftarrow c}$ need not equal
$\Gamma_{c\leftarrow a}^{-1}$.  Therefore the multiplicative relative defect
$D^\times_{c,b,a}$ and triangle holonomy $H_{c,b,a}$ are different
comparisons.  The formal module proves that global flatness trivializes both.
It does not identify them for an arbitrary connection.
\end{boundary}

\begin{figure}[H]
\centering
\includegraphics[width=.96\linewidth]{transport_vs_connection.pdf}
\caption{Exact coordinate transport and connection transport are distinct.
The canonical loop telescopes to the identity.  A chosen connection may
carry nontrivial holonomy, which transforms by conjugation under gauge.}
\label{fig:transport-connection}
\end{figure}

\Cref{fig:transport-connection} records this separation visually: the exact
coordinate route telescopes, whereas the chosen connection route may carry a
nonidentity loop.

\begin{keyidea}[title={Coordinate transport and connection choice}]
Canonical transport compares exact coordinate descriptions of the same
object.  Its loops telescope.  A connection supplies a separate family of
comparison maps, so its loops can carry information about that choice.
The distinction keeps coordinate identities separate from connection
assumptions.
\end{keyidea}

\subsection{A stronger reversible pair API}

Let $R$ be a commutative ring, let $B$ be a type of charts or refinement
states, and let each $C_a$ be an $R$-module.  The chart type is not assumed
finite.  The following definition is deliberately not derived from analysis
and synthesis.

\begin{definition}[Normalized reversible connection]
\label{def:connection}
A normalized reversible connection is a family of linear equivalences
\[
\Gamma_{b\leftarrow a}:C_a\xrightarrow{\;\sim\;}C_b
\qquad(a,b\in B)
\]
such that
\[
\Gamma_{a\leftarrow a}=\id,
\qquad
\Gamma_{a\leftarrow b}=\Gamma_{b\leftarrow a}^{-1}.
\]
No cocycle equation is included in the structure.
\end{definition}
\statusleanthree{Komorebi.PaperIII.Transport.ReversibleConnection}

This adds heterogeneous fibers and inversion on reversed edges to the
homogeneous minimal connection.  The connection is still a separate choice
of maps, rather than a definition of canonical coordinate transport.

\begin{definition}[Calibration and flatness]
\label{def:calibration-flatness}
When the $C_a$ are exact packet coordinates of one ambient module, the
connection is \emph{canonically calibrated} if
\[
\Gamma_{b\leftarrow a}=J_{b\leftarrow a}
\quad\text{for every }a,b.
\]
It is \emph{flat} if
\[
\Gamma_{c\leftarrow b}\Gamma_{b\leftarrow a}
=
\Gamma_{c\leftarrow a}
\quad\text{for every }a,b,c.
\]
The canonical connection is flat.  A general flat connection is proved in
\cref{thm:flat-gauge-canonical} to be related to it by a gauge change.
\end{definition}
\statusleanthree{Komorebi.PaperIII.Transport.ReversibleConnection.CocycleAt; Komorebi.PaperIII.Transport.ReversibleConnection.Flat; Komorebi.PaperIII.Transport.ReversibleConnection.canonicalConnection; Komorebi.PaperIII.Transport.ReversibleConnection.canonicalConnection_flat}

The words \emph{calibrated} and \emph{flat} have only the algebraic meanings
displayed above.

\begin{remark}[A connection on a chart graph]
One may specify $\Gamma_e$ only on the oriented edges of a chart graph,
together with inverse maps on reversed edges.  Path transport and loop
holonomy below still make sense.  A direct cocycle defect for $(a,b,c)$,
however, requires a declared direct transport from $a$ to $c$.  The pairwise
API in \cref{def:connection} keeps this comparison total.
\end{remark}

\subsection{Cocycle defect and triangle holonomy}

The commutative ring hypothesis makes subtraction of linear maps available
from this point through \cref{sec:obstruction-packets}.

\begin{definition}[Cocycle defect]
\label{def:cocycle-defect}
For $a,b,c\in B$, define the additive cocycle defect
\[
\Delta_{c,b,a}
=
\Gamma_{c\leftarrow b}\Gamma_{b\leftarrow a}
-\Gamma_{c\leftarrow a}
\in\Hom_R(C_a,C_c).
\]
\end{definition}
\statusleanthree{Komorebi.PaperIII.Transport.ReversibleConnection.cocycleDefect}

The manuscript subscript order is target, intermediate, source.  The Lean
arguments are supplied in traversal order \leanname{(a,b,c)}.

\begin{definition}[Triangle holonomy]
\label{def:triangle-holonomy}
The triangle holonomy based at $a$ is
\[
H_{c,b,a}
=
\Gamma_{a\leftarrow c}
\Gamma_{c\leftarrow b}
\Gamma_{b\leftarrow a}
\in\operatorname{Aut}_R(C_a).
\]
It traverses $a\to b\to c\to a$.
\end{definition}
\statusleanthree{Komorebi.PaperIII.Transport.ReversibleConnection.triangleHolonomy; Komorebi.PaperIII.Transport.ReversibleConnection.triangleObstruction}

Holonomy is attached to the chosen connection $\Gamma$, never to an
unspecified coordinate change.

\begin{theorem}[Identity between defect and holonomy]
\label{thm:defect-holonomy}
For every triple $(a,b,c)$,
\[
\boxed{
H_{c,b,a}-\id_{C_a}
=
\Gamma_{a\leftarrow c}\Delta_{c,b,a}.
}
\]
Consequently,
\[
\Delta_{c,b,a}=0
\quad\Longleftrightarrow\quad
H_{c,b,a}=\id_{C_a}.
\]
\end{theorem}
\statusleanthree{Komorebi.PaperIII.Transport.ReversibleConnection.triangleHolonomy_sub_id; Komorebi.PaperIII.Transport.ReversibleConnection.cocycleDefect_eq_zero_iff_triangleHolonomy_eq_refl}

\begin{proof}
Using reversibility,
\[
\begin{aligned}
\Gamma_{a\leftarrow c}\Delta_{c,b,a}
&=
\Gamma_{a\leftarrow c}
\Gamma_{c\leftarrow b}
\Gamma_{b\leftarrow a}
-
\Gamma_{a\leftarrow c}\Gamma_{c\leftarrow a}\\
&=H_{c,b,a}-\id_{C_a}.
\end{aligned}
\]
Left composition by the equivalence $\Gamma_{a\leftarrow c}$ is injective,
which gives the equivalence of the two vanishing statements.
\end{proof}

\begin{corollary}[Canonical transport has zero connection defect]
\label{cor:canonical-zero-defect}
If $\Gamma_{b\leftarrow a}=J_{b\leftarrow a}$, then
\[
\Delta_{c,b,a}=0,
\qquad
H_{c,b,a}=\id
\]
for all triples.
\end{corollary}
\statusleanthree{Komorebi.PaperIII.Transport.ReversibleConnection.canonicalTransport_zero_connectionDefect; Komorebi.PaperIII.Transport.ReversibleConnection.canonicalTransport_triangleHolonomy_eq_refl}

\begin{proof}
The defect vanishes by the canonical cocycle in
\cref{thm:J-cocycle}. Apply \cref{thm:defect-holonomy}.
\end{proof}

\subsection{Paths and endpoint independence}

\begin{definition}[Connection transport along a path]
\label{def:path-transport}
For a vertex sequence
$p=(a_0,a_1,\ldots,a_n)$, define
\[
\Gamma_p
=
\Gamma_{a_n\leftarrow a_{n-1}}\cdots
\Gamma_{a_1\leftarrow a_0},
\qquad
\Gamma_{(a_0)}=\id_{C_{a_0}}.
\]
If $a_n=a_0$, then $\Gamma_p$ is the holonomy of the loop $p$.
\end{definition}
\statusleanthree{Komorebi.PaperIII.Transport.ReversibleConnection.pathTransport; Komorebi.PaperIII.Transport.ReversibleConnection.VertexPath; Komorebi.PaperIII.Transport.ReversibleConnection.indexedPathTransport}

Products follow path order: the first traversed edge acts on the right.

\begin{proposition}[Flat pair transport collapses to endpoints]
\label{prop:flat-endpoints}
If $\Gamma$ is flat, then every path satisfies
\[
\Gamma_p=\Gamma_{a_n\leftarrow a_0}.
\]
In particular, every loop has identity holonomy.
\end{proposition}
\statusleanthree{Komorebi.PaperIII.Transport.ReversibleConnection.pathTransport_eq_direct_of_flat; Komorebi.PaperIII.Transport.ReversibleConnection.indexedPathTransport_eq_direct_of_flat; Komorebi.PaperIII.Transport.ReversibleConnection.indexedLoopTransport_eq_refl_of_flat}

\begin{proof}
Induct on the path length.  Appending $a_{n+1}$ left-multiplies the
inductive expression by $\Gamma_{a_{n+1}\leftarrow a_n}$, and flatness
replaces the two factors by
$\Gamma_{a_{n+1}\leftarrow a_0}$.  For a loop, normalization gives
$\Gamma_{a_0\leftarrow a_0}=\id$.
\end{proof}

\begin{remark}
The converse \emph{trivial holonomy implies flatness} depends on which loops
are present in an API defined on a chart graph.  In the complete pair API,
triviality of every triangle holonomy is equivalent to flatness by
\cref{thm:defect-holonomy}.  No theorem about topological simple connectivity is
being invoked.
\end{remark}

\section{Gauge covariance and comparison with canonical transport}
\label{sec:gauge}

\begin{definition}[Gauge change]
\label{def:gauge-change}
A gauge is a family of automorphisms $G_a\in\operatorname{Aut}_R(C_a)$.
It acts on a connection by
\[
\Gamma^G_{b\leftarrow a}
=
G_b\Gamma_{b\leftarrow a}G_a^{-1}.
\]
\end{definition}
\statusleanthree{Komorebi.PaperIII.Transport.ReversibleConnection.Gauge; Komorebi.PaperIII.Transport.ReversibleConnection.gauge}

A gauge changes coordinate representatives at each chart.  The chart graph
stays the same.

\begin{theorem}[Gauge covariance]
\label{thm:gauge-covariance}
The defect and triangle holonomy transform as
\[
\Delta^G_{c,b,a}
=G_c\Delta_{c,b,a}G_a^{-1},
\qquad
H^G_{c,b,a}
=G_aH_{c,b,a}G_a^{-1}.
\]
Thus flatness and triviality of triangle holonomy are gauge invariant.
\end{theorem}
\statusleanthree{Komorebi.PaperIII.Transport.ReversibleConnection.gauge_cocycleDefect; Komorebi.PaperIII.Transport.ReversibleConnection.gauge_triangleHolonomy; Komorebi.PaperIII.Transport.ReversibleConnection.gauge_flat_iff; Komorebi.PaperIII.Transport.ReversibleConnection.gauge_triangleHolonomy_eq_refl_iff}

\begin{proof}
Substitute the definition of $\Gamma^G$.  Adjacent factors
$G_b^{-1}G_b$ and $G_c^{-1}G_c$ cancel.  In the defect, both the composite
and direct terms acquire the same outer factors $G_c$ and $G_a^{-1}$.  In
the loop, all inner factors cancel and only conjugation at the base remains.
\end{proof}

\begin{theorem}[Flat connections are gauges of the canonical pair transport]
\label{thm:flat-gauge-canonical}
Suppose the $C_a$ are exact packet coordinate modules for one ambient module,
so the canonical $J_{b\leftarrow a}$ is available.  If $\Gamma$ is a flat
normalized reversible pair connection, choose $o\in B$ and set
\[
G_a=\Gamma_{a\leftarrow o}J_{o\leftarrow a}.
\]
Then every pair transport satisfies
\[
\boxed{
\Gamma_{b\leftarrow a}
=G_bJ_{b\leftarrow a}G_a^{-1}.
}
\]
\end{theorem}
\statusleanthree{Komorebi.PaperIII.Transport.ReversibleConnection.flatGauge; Komorebi.PaperIII.Transport.ReversibleConnection.flat_isGaugeOfCanonical}

\begin{proof}
Reversibility gives
$G_a^{-1}=J_{a\leftarrow o}\Gamma_{o\leftarrow a}$.  Therefore
\[
\begin{aligned}
G_bJ_{b\leftarrow a}G_a^{-1}
&=
\Gamma_{b\leftarrow o}
J_{o\leftarrow b}J_{b\leftarrow a}
J_{a\leftarrow o}
\Gamma_{o\leftarrow a}\\
&=
\Gamma_{b\leftarrow o}\Gamma_{o\leftarrow a}\\
&=
\Gamma_{b\leftarrow a}.
\end{aligned}
\]
The first reduction uses the canonical cocycle twice. The second uses
flatness of $\Gamma$.
\end{proof}

\begin{boundary}[Flat does not mean calibrated]
The theorem produces a gauge relating $\Gamma$ to $J$.  It does not prove
$\Gamma=J$ before that gauge is chosen.  Conversely, nontrivial triangle
holonomy cannot be removed by a gauge, because conjugation cannot turn a
nonidentity automorphism into the identity.
\end{boundary}

\section{Obstruction packets: local certificates and holonomy blocks}
\label{sec:obstruction-packets}

\subsection{The generic compiled certificate}

Let $\mathsf{Candidate}$, $\mathsf{Packet}$, and $\mathsf{Residual}$ be
arbitrary types, with a distinguished
$0\in\mathsf{Residual}$, and let
\[
r:\mathsf{Candidate}\longrightarrow
\mathsf{Packet}\longrightarrow\mathsf{Residual}
\]
be a residual function.

\begin{definition}[Solution and obstruction packet for a candidate]
\label{def:generic-obstruction}
A candidate $x$ is a solution when
\[
\operatorname{IsSolution}_r(x)
\quad\Longleftrightarrow\quad
\forall p,\ r(x,p)=0.
\]
An obstruction packet for $x$ is a pair consisting of a named packet $p$
and a certificate $r(x,p)\ne0$.
\end{definition}
\statusleanthree{Komorebi.PaperIII.NonlinearObstructions.IsSolution; Komorebi.PaperIII.NonlinearObstructions.ObstructionPacket}

\begin{theorem}[Exact logical force of a local obstruction]
\label{thm:generic-obstruction}
Every obstruction packet for $x$ proves that $x$ is not a solution, and
\[
\neg\operatorname{IsSolution}_r(x)
\quad\Longleftrightarrow\quad
\text{there exists an obstruction packet for }x.
\]
The reverse implication from failure to a packet witness is classical and
does not require the packet type to be finite.
\end{theorem}
\statusleanthree{Komorebi.PaperIII.NonlinearObstructions.ObstructionPacket.not_isSolution; Komorebi.PaperIII.NonlinearObstructions.obstructionPacket_of_not_isSolution; Komorebi.PaperIII.NonlinearObstructions.not_isSolution_iff_nonempty_obstructionPacket}

\begin{proof}
A purported solution would make the residual vanish at the certified packet,
contradicting its nonzero field.  Conversely, if no packet carries a nonzero
residual, then every residual is zero and the candidate satisfies the
definition of solution.
\end{proof}

\begin{interpretation}[title={What one obstruction certifies}]
The certificate answers one question: does this specified candidate satisfy
every declared packet condition?  One nonzero residual settles that question
for this candidate.  A statement that no candidate can solve the problem
would require an obstruction for every candidate.
\end{interpretation}

\begin{boundary}[Candidate-specific means candidate-specific]
The theorem rules out the named $x$.  It neither constructs a solution nor
shows that every candidate is obstructed.  It provides no minimality,
stability, computability, or classification theorem for the packet witness.
\end{boundary}

\subsection{Specialization to holonomy blocks}

For a reversible connection, specialize the generic residual at a fixed
triangle by taking a packet label to be a pair $(\mu,\nu)$ and the residual
to be the corresponding block of $H-\id$.  A nonzero block, together with its
label and a proof of nonvanishing, gives an obstruction packet certificate
in the sense of \cref{def:generic-obstruction}.  The reconstruction theorem
below supplies the further step: vanishing of all these local residuals is
equivalent to the cocycle equation at that triangle.

Fix a triangle $(a,b,c)$ and resolve the base fiber $C_a$ by a finite complete
family of idempotents $(\Pi^a_\mu)_{\mu\in I_a}$.  Orthogonality may also hold,
as it does for the exact range coordinate projectors below, but the compiled
reconstruction theorem does not require it.  In exact range coordinates,
$\Pi^a_\mu$ is simply the projector that retains component $\mu$ and zeros
the other components.

\begin{definition}[Holonomy obstruction and packet residuals]
\label{def:obstruction-packet}
Define
\[
\Ob_{c,b,a}=H_{c,b,a}-\id_{C_a},
\qquad
\Ob^{\mu\nu}_{c,b,a}
=
\Pi^a_\mu\Ob_{c,b,a}\Pi^a_\nu.
\]
The map $\Ob^{\mu\nu}_{c,b,a}$ is the obstruction block, or packet residual,
from input packet $\nu$ to output packet $\mu$.  If it is nonzero, then
$(\mu,\nu)$ and that nonvanishing proof form an obstruction packet
certificate.  Diagonal blocks detect holonomy within one packet.  Off-diagonal
blocks detect packet mixing induced by holonomy.
\end{definition}
\statusleanthree{Komorebi.PaperIII.HolonomyBlocks.PacketResolution; Komorebi.PaperIII.HolonomyBlocks.PacketResolution.block}

The certificate records a specified holonomy defect, the label of an input
and output packet, and a proof that the corresponding block is nonzero.

\begin{theorem}[Exact reconstruction and detection by obstruction blocks]
\label{thm:obstruction-detection}
The total obstruction has the exact finite resolution
\[
\boxed{
\Ob_{c,b,a}
=
\sum_{\mu,\nu\in I_a}
\Ob^{\mu\nu}_{c,b,a}.
}
\]
Moreover, the following are equivalent:
\begin{enumerate}[label=(\roman*)]
\item $\Delta_{c,b,a}=0$;
\item $H_{c,b,a}=\id$;
\item $\Ob_{c,b,a}=0$;
\item $\Ob^{\mu\nu}_{c,b,a}=0$ for every $\mu,\nu$.
\end{enumerate}
\end{theorem}
\statusleanthree{Komorebi.PaperIII.HolonomyBlocks.obstruction_sum_blocks; Komorebi.PaperIII.HolonomyBlocks.obstructionBlocks_zero_iff; Komorebi.PaperIII.HolonomyBlocks.obstructionDetection_equivalences}

\begin{proof}
Insert completeness on both sides:
\[
\Ob
=
\left(\sum_\mu\Pi^a_\mu\right)
\Ob
\left(\sum_\nu\Pi^a_\nu\right)
=
\sum_{\mu,\nu}\Pi^a_\mu\Ob\Pi^a_\nu.
\]
Thus (iv) implies (iii), and (iii) implies (iv) by composition with the
projectors.  Conditions (ii) and (iii) are equivalent by the definition of
$\Ob$ and additive cancellation.
Conditions (i) and (ii) are equivalent by
\cref{thm:defect-holonomy}.
\end{proof}

\begin{keyidea}[title={Why both packet indices matter}]
For the declared finite complete resolution, reconstruction uses both
packet indices.  Inspecting only diagonal blocks can miss transport from one
packet to another.  The full family records that mixing and reconstructs the
total obstruction exactly.
\end{keyidea}

\Cref{fig:holonomy-obstruction-blocks} depicts the finite block resolution
on both sides of the obstruction and the off-diagonal blocks that witness
packet mixing.

\begin{figure}[H]
\centering
\includegraphics[width=.96\linewidth]{holonomy_obstruction_blocks.pdf}
\caption{A finite complete packet resolution reconstructs the total holonomy
obstruction from blocks indexed by input and output packets.  A nonzero
off-diagonal block is an exact witness of packet mixing.}
\label{fig:holonomy-obstruction-blocks}
\end{figure}

\begin{corollary}[Witness supported in one packet]
\label{cor:packet-witness}
If $\Ob^{\mu\nu}_{c,b,a}\ne0$, then there is a vector $x$ supported in packet
$\nu$ such that
\[
\Pi^a_\mu\Ob_{c,b,a}x\ne0.
\]
Such an $x$ is an exact finite witness of nonflatness.
\end{corollary}
\statusleanthree{Komorebi.PaperIII.HolonomyBlocks.exists_packetWitness_of_obstructionBlock_ne_zero}

\begin{proof}
Classically choose $y$ on which the nonzero linear map
$\Pi^a_\mu\Ob\Pi^a_\nu$ is nonzero and set $x=\Pi^a_\nu y$.
Idempotence gives $\Pi^a_\nu x=x$, and the required output is the chosen
nonzero value.
\end{proof}

\subsection{Gauge behavior of packet blocks}

Vanishing or nonvanishing of the total triangle obstruction is gauge invariant
by \cref{thm:gauge-covariance}, but a fixed packet pattern need not be.  A
general $G_a$ can mix the projectors $\Pi^a_\mu$.

\begin{proposition}[Gauge covariance when packet projectors are preserved]
\label{prop:packet-gauge}
If $G_a\Pi^a_\mu=\Pi^a_\mu G_a$ for every $\mu$, then
\[
\left(\Ob^G\right)^{\mu\nu}
=
G_a\Ob^{\mu\nu}G_a^{-1}.
\]
Consequently the zero/nonzero status of each obstruction block is preserved
by gauges that preserve the packet projectors.
\end{proposition}
\statusleanthree{Komorebi.PaperIII.HolonomyBlocks.obstructionBlock_gauge_of_commutes; Komorebi.PaperIII.HolonomyBlocks.obstructionBlock_gauge_eq_zero_iff}

\begin{proof}
Use $\Ob^G=G_a\Ob G_a^{-1}$ and commute $G_a$ and $G_a^{-1}$ through the
two packet projectors.  Composition by an equivalence is injective in both
directions.
\end{proof}

\begin{boundary}[Packet patterns are not absolute]
Without the commutation hypothesis in \cref{prop:packet-gauge}, only the
vanishing of the total obstruction is gauge invariant.  The location of
individual nonzero blocks may change under a gauge that mixes packets.
\end{boundary}

\section{Products supported in packets and holonomy walks}
\label{sec:holonomy-walks}

The passage from words to walks also applies to products of different matrices.
Let $V$
be one finite packet type, let $R$ be a semiring, and let
$B_1,\ldots,B_m\in R^{V\times V}$ be a finite sequence of possibly different
packet matrices.  Rows are targets and columns are sources as before.

\begin{proposition}[Packet path expansion for different matrices]
\label{prop:heterogeneous-path}
For $i,j\in V$ and $m\geq1$, set $i_0=i$ and $i_m=j$.  Then
\[
\boxed{
(B_m\cdots B_1)_{ji}
=
\sum_{(i_1,\ldots,i_{m-1})\in V^{m-1}}
  B_m[i_m,i_{m-1}]\cdots B_1[i_1,i_0].
}
\]
When $m=1$, the sum is over the unique empty tuple and its summand is
$B_1[j,i]$.
For $m=0$, the product is the identity and the entry is the Kronecker delta
$\delta_{ji}$.  If every $B_r[q,p]$ vanishes outside a declared relation
$\mathcal R_r(p,q)$, then the endpoint entry is zero whenever there is no
packet path compatible with the successive relations.
\end{proposition}
\statusleanthree{Komorebi.PaperIII.EdgeOperators.heterogeneousProduct; Komorebi.PaperIII.EdgeOperators.heterogeneousPathSum; Komorebi.PaperIII.EdgeOperators.heterogeneousPacketPathExpansion; Komorebi.PaperIII.EdgeOperators.SupportedStep; Komorebi.PaperIII.EdgeOperators.SupportedPath; Komorebi.PaperIII.EdgeOperators.heterogeneousProduct_apply_eq_zero_of_no_supportedPath}

\begin{proof}
Expand each matrix product entry and distribute the finite sums.  Induction on
the factor list yields the displayed path sum, with the empty list evaluated
as the identity matrix.  If no supported intermediate packet exists, every
summand contains an entry that is zero by its declared support relation.
\end{proof}

\begin{boundary}[Changing packet types along a chain]
The compiled theorems
\leanname{heterogeneousPacketPathExpansion_two} and
\leanname{heterogeneousPacketPathExpansion_three} allow different packet
types at the successive stages of two- and three-factor products.  The
arbitrary-length theorem above permits different factors but uses one common
finite packet type $V$.  A fully dependent arbitrary-length chain of changing
modules and packet types remains outside the finite algebraic contract.
\end{boundary}

\begin{remark}[Checked specialization]
When every $B_r$ is one repeated endomorphism $T$, the recursive
specialization to packet powers is inherited from Chapter I.
\end{remark}
\statusleanprior{Komorebi.PacketNaturalityCompletion.Packets.transferBlock_pow_eq_packetPathSum}

For a loop represented on one common finite scalar packet coordinate type,
the proposition resolves its successive transport matrices into packet
walks.  If $R$ is a ring, as in the reversible connection API, then
$(H-I)_{ji}=H_{ji}-\delta_{ji}$.  In particular, off-diagonal entries are
unchanged.  Thus a nonzero off-diagonal holonomy obstruction in this model
with a common index type requires a transport word along the closed chart
loop whose packet path runs from its input packet to its output packet.  This
conclusion requires declared block supports for the connection matrices.
An arbitrary connection need not respect the transfer graph from
\cref{sec:edge-operators}.

\section{A minimal exact example}
\label{sec:example}

\begin{example}[A scalar triangle]
Let $B=\{0,1,2\}$ and $C_0=C_1=C_2=\bbk$ for a field $\bbk$.  Choose
$u,v,w\in\bbk^\times$ and set
\[
\Gamma_{1\leftarrow0}(x)=ux,\qquad
\Gamma_{2\leftarrow1}(x)=vx,\qquad
\Gamma_{2\leftarrow0}(x)=wx,
\]
with reverse transports given by inverse scalars.  Then
\[
\Delta_{2,1,0}(x)=(vu-w)x,
\]
and the triangle holonomy based at $0$ is
\[
H_{2,1,0}(x)=w^{-1}vux.
\]
Therefore
\[
H_{2,1,0}-\id
=w^{-1}\Delta_{2,1,0}.
\]
The connection is flat on this triangle exactly when $vu=w$.  If the three
coordinate systems are literally the same and canonical transport is used,
then $u=v=w=1$ and the obstruction vanishes.  Choosing arbitrary nonzero
$u,v,w$ is a connection choice, not a consequence of changing exact packet
coordinates.
\end{example}

\begin{example}[A transfer walk with two steps]
Let a packet graph have edges $e:0\to1$ and $f:1\to2$, with no other edges.
Then
\[
T=L_e+L_f,\qquad
(T^2)_{20}=c(f)c(e).
\]
The products $L_eL_e$, $L_eL_f$, and $L_fL_f$ vanish by typed mismatch.
Only $L_fL_e$, corresponding to the walk $0\to1\to2$, contributes to the
$(2,0)$ entry.  This example also fixes the noncommutative coefficient order.
\end{example}

\section{Canonical coordinates for arbitrary maps}
\label{sec:arbitrary-maps}

The analysis and synthesis identities are identities of functions, not only
of linear maps.  Consequently exact coordinate change does not become
curved merely because the represented operation is nonlinear.

\begin{proposition}[Exact naturality for arbitrary functions]
\label{prop:set-naturality}
Let $\cA^X_a:X\xrightarrow{\sim}X_a$ and
$\cA^X_b:X\xrightarrow{\sim}X_b$ be exact coordinate equivalences for $X$,
with inverses $\cS^X_a$ and $\cS^X_b$.  Likewise, let
$\cA^Y_r:Y\xrightarrow{\sim}Y_r$ and
$\cA^Y_s:Y\xrightarrow{\sim}Y_s$ be exact coordinate equivalences for $Y$,
with inverses $\cS^Y_r$ and $\cS^Y_s$.  Let $F:X\to Y$ be any function, with
no linearity or differentiability hypothesis.  Set
\[
J^X_{b\leftarrow a}=\cA^X_b\circ\cS^X_a,
\qquad
J^Y_{s\leftarrow r}=\cA^Y_s\circ\cS^Y_r,
\]
and define
\[
\widehat F_{r\leftarrow a}
=
\cA^Y_r\circ F\circ\cS^X_a.
\]
Then
\[
J^Y_{s\leftarrow r}\circ\widehat F_{r\leftarrow a}
=
\widehat F_{s\leftarrow b}\circ J^X_{b\leftarrow a}.
\]
If $\cA^Z_t:Z\xrightarrow{\sim}Z_t$ is another exact coordinate equivalence
and $G:Y\to Z$ is any function, then exact coordinate representation also
preserves composition:
\[
\widehat{G\circ F}_{t\leftarrow a}
=
\widehat G_{t\leftarrow r}\circ\widehat F_{r\leftarrow a}.
\]
\end{proposition}
\statusleanthree{Komorebi.PaperIII.Transport.exactCoordinateMap; Komorebi.PaperIII.Transport.exactCoordinateMap_comp; Komorebi.PaperIII.Transport.exactCoordinateMap_naturality}

\begin{proof}
Both sides of the naturality equation reduce to
$\cA^Y_s\circ F\circ\cS^X_a$ using
$\cS\circ\cA=\id$ in $X$ and $Y$.  For composition, the middle
$\cS^Y\circ\cA^Y$ cancels in the same way.
\end{proof}

\begin{boundary}[Where nonlinear defects can enter]
A nonlinear \emph{coordinate} defect cannot be blamed on the exact canonical
$J$ of \cref{def:canonical-J}. By \cref{prop:set-naturality}, that defect is
identically zero.  Nonzero defects require extra structure such as a
noncanonical connection, truncation, approximation, restricted domains, or
failure of an operation to respect a declared packet model.
\end{boundary}

\clearpage\section{Nonlinear operation trees: locality and the formalization frontier}
\label{sec:operation-trees}

\subsection{The compiled foundation of unary and binary trees}

\begin{definition}[Finite grammar for operation trees]
\label{def:compiled-operation-tree}
For types of input labels $I$, constants $K$, unary symbols $U$, and binary
symbols $B$, the finite formal grammar has four constructors:
\[
\begin{aligned}
t ::={}&\operatorname{input}(i)
\mid\operatorname{constant}(k)
\mid\operatorname{unary}(u,t)\\
&\mid\operatorname{binary}(b,t_1,t_2).
\end{aligned}
\]
Given an arbitrary value type $Y$, evaluation takes completely explicit maps
\[
I\to Y,\qquad K\to Y,\qquad U\to Y\to Y,\qquad
B\to Y\to Y\to Y.
\]
No algebraic, continuity, or differentiability law is imposed on these
semantics.
\end{definition}
\statusleanthree{Komorebi.PaperIII.NonlinearObstructions.OperationTree; Komorebi.PaperIII.NonlinearObstructions.OperationTree.eval}

When $I$ has decidable equality, let $\operatorname{inputs}(t)$ be the finite
set of input labels actually occurring in $t$.

\begin{theorem}[Input locality for operation trees]
\label{thm:tree-locality}
Let $\eta_1,\eta_2:I\to Y$ be two input environments.  If
\[
\eta_1(i)=\eta_2(i)
\qquad\text{for every }i\in\operatorname{inputs}(t),
\]
then the two evaluations of $t$ are equal under the same constant, unary, and
binary semantics.  In particular, if
$\operatorname{inputs}(t)=\varnothing$, evaluation is independent of the
input environment.
\end{theorem}
\statusleanthree{Komorebi.PaperIII.NonlinearObstructions.OperationTree.eval_congr_on_inputs; Komorebi.PaperIII.NonlinearObstructions.OperationTree.eval_eq_of_inputs_eq_empty}

\begin{proof}
Structural induction on $t$.  The input case is the agreement hypothesis.
The constant case is reflexive. The unary case applies the induction
hypothesis to its child. The binary case applies it independently to the
left and right subtrees, whose input sets form a union.
\end{proof}

\begin{boundary}[Exact force of locality]
The theorem says evaluation reads only syntactically named inputs.  It does
not prove that the operations are linear or nonlinear in any analytic sense,
that the tree is stable, that two different trees are equivalent, or that a
solution to a tree equation exists.
\end{boundary}

\subsection{Typed colored extension beyond the finite contract}

\begin{definition}[Finitary colored operation signature]
\label{def:signature}
A finitary colored signature consists of:
\begin{itemize}
\item a finite type $V$ of colors;
\item a type $\Sigma$ of operation symbols;
\item an arity $r(\sigma)\in\N$ for every $\sigma\in\Sigma$;
\item input colors
$\operatorname{in}_\sigma:\operatorname{Fin}(r(\sigma))\to V$;
\item an output color $\operatorname{out}(\sigma)\in V$.
\end{itemize}
A realization assigns a space $X_v$ to each color and a map
\[
\Phi_\sigma:
\prod_{q\in\operatorname{Fin}(r(\sigma))}X_{\operatorname{in}_\sigma(q)}
\longrightarrow X_{\operatorname{out}(\sigma)}.
\]
\end{definition}
\statusroadmap{The intended Lean target is a dependent inductive typed tree. No such declaration is part of the finite algebraic completion contract.}

\begin{definition}[Typed operation tree and evaluation]
\label{def:operation-tree}
A typed operation tree is recursively either a leaf of declared color or a
node labeled by $\sigma$, with one correctly colored child for each input
slot of $\sigma$.  Evaluation is recursive:
\[
\ev(\operatorname{node}(\sigma;t_1,\ldots,t_{r(\sigma)}))
=
\Phi_\sigma(\ev(t_1),\ldots,\ev(t_{r(\sigma)})).
\]
\end{definition}
\statusroadmap{Candidate API: \leanname{Komorebi.PaperIII.NonlinearObstructions.ColoredOperationTree} and \leanname{ColoredOperationTree.eval}.  This is stronger than the compiled uncolored unary--binary grammar.}

\subsection{Packet labelings and the multilinear first target}

Suppose every $X_v$ has a finite complete packet family
$(P^v_i)_{i\in I_v}$.  A packet labeling assigns an admissible packet index
to each edge or node of a typed tree.  For multilinear $\Phi_\sigma$, inserting
\[
\id_{X_v}=\sum_{i\in I_v}P^v_i
\]
at every input and output yields a finite sum over trees with packet labels.
This is the branching analogue of inserting intermediate projectors in
\cref{prop:heterogeneous-path}.

\begin{roadmap}[Finite multilinear tree expansion]
\label{roadmap:multilinear-tree}
The first milestone for the nonlinear module is an exact theorem saying that the
evaluation of a finite tree of multilinear nodes equals the finite sum of its
packet-resolved labeled evaluations.  Required proof obligations include:
\begin{enumerate}[label=(\alph*)]
\item a dependent well-typed tree syntax;
\item finite packet-label types at every color;
\item an explicit convention for nullary and unary nodes;
\item multilinearity in every input slot;
\item preservation of the order of input slots without assuming commutativity of
operations.
\end{enumerate}
\end{roadmap}
\statusroadmap{Candidate target: \leanname{Komorebi.PaperIII.NonlinearObstructions.ColoredOperationTree.eval_eq_sum_packetLabellings}.}

\subsection{Connection defects at nonlinear nodes}

For coordinate models of the source and target colors, an additive target
coordinate space, and a chosen noncanonical connection, a pointwise node
defect can be defined by
\[
\mathfrak D^\Gamma_{b,a}(\Phi)
=
\Gamma^Y_{b\leftarrow a}\circ\widehat\Phi_a
-
\widehat\Phi_b\circ
\left(\prod_q\Gamma^{X_q}_{b\leftarrow a}\right).
\]
This expression compares two \emph{chosen} routes.  With canonical exact
transport in every slot it vanishes by \cref{prop:set-naturality}.

\begin{roadmap}[Operation-tree defect propagation]
\label{roadmap:tree-defect}
For multilinear nodes, one expects an exact finite expansion obtained by
replacing one or more nodes by their defects and transporting the remaining
branches.  For smooth nonlinear nodes, the corresponding chain rule is
indexed by rooted partitions and Fa\`a di Bruno trees.  A formal theorem will
require a fixed differentiability category, base points, derivative
conventions, symmetry factors, and remainder policy.
\end{roadmap}
\statusroadmap{Tree-level defect identities, higher chain rules, and nonlinear norm bounds are outside the finite Lean contract. Ordinary mathematical results in these directions appear in Part~II under their stated hypotheses.}

\begin{boundary}[Analytic and geometric exclusions]
The finite formal contract does not include convergence of an infinite tree series, existence
of derivatives, a Banach-space Fa\`a di Bruno theorem, a Baker--Campbell--Hausdorff
expansion, a continuum limit, or an identification of algebraic holonomy with
curvature of a smooth bundle. Part~II supplies ordinary proofs for its explicit nonlinear and analytic statements. The theory does not derive a connection from
data, optimization, physics, or probability.
\end{boundary}

\section{Finite algebraic formalization ledger}
\label{sec:ledger}

\small
\begin{longtable}{L{0.17\textwidth}L{0.18\textwidth}L{0.25\textwidth}L{0.22\textwidth}}
\toprule
Topic & Status & Declaration or roadmap target & Exact boundary\\
\midrule
\endfirsthead
\toprule
Topic & Status & Declaration or roadmap target & Exact boundary\\
\midrule
\endhead
Packet matrix and edge operator
& Lean-backed III
& \leanname{PacketGraph.packetMatrix}; \leanname{EdgeSystem.edgeOp}
& Finite scalar matrices over a semiring, with one scalar coordinate per packet.\\
\addlinespace
Edge support and pair dichotomy
& Lean-backed III
& \leanname{edgeOp_eq_zero_off_support}; \leanname{edgeOp_mul_edgeOp_eq_zero}; \leanname{edgeOp_mul_edgeOp_of_composable}
& Broken interfaces vanish. Matched pairs collapse to the endpoint block with coefficient in operator order.\\
\addlinespace
Typed word evaluation
& Lean-backed III
& \leanname{Composable}; \leanname{wordOp}; \leanname{wordOp_append}; \leanname{wordOp_eq_zero_of_not_composable}
& List order is execution order. Concatenation reverses operator-product order. Nonwalks vanish.\\
\addlinespace
Finite word and walk expansions
& Lean-backed III
& \leanname{explicitWordExpansion}; \leanname{walkExpansion}; \leanname{wordExpansion_eq_explicitWordExpansion}; \leanname{transferOp_pow_eq_walkExpansion}
& Actual finite tuple sum and zero-filtered composable-walk sum over finite graph and packet types.\\
\addlinespace
Endpoint coefficient walk sum
& Lean-backed III
& \leanname{reverseCoefficient}; \leanname{Connects}; \leanname{transferOp_pow_apply_eq_sum_walks}
& Exact scalar endpoint entry, including the empty-walk case and noncommutative coefficient order.\\
\addlinespace
Fixed reachability and orbit support
& Lean-backed III, with an inherited hull bridge
& \leanname{transferOp_pow_apply_eq_zero_of_no_walk}; \leanname{transferOp_pow_mulVec_apply_eq_zero_of_unreachable_support}
& Finite powers and supported seeds. The orbit-hull span conclusion uses the frozen Chapter I contract.\\
\addlinespace
Exact frame transport
& Lean-backed III
& \leanname{exactTransport_cocycle}; \leanname{exactPathTransport_eq_endpoint}; \leanname{exactPathTransport_loop}
& One frame family on a single module. Arbitrary finite paths telescope and closed paths are trivial.\\
\addlinespace
Exact coordinate cocycle
& Inherited Lean-backed II
& \leanname{NestedRefinement.refineCoordinates_comp}; \leanname{NestedRefinement.coordinateRefinementEquiv_trans}
& Complete finite projector coordinates, not reduction.\\
\addlinespace
Typed coordinate naturality
& Inherited Lean-backed II
& \leanname{NestedRefinement.coordinateOperator_comp}; \leanname{NestedRefinement.coordinateOperator_naturality}
& Linear maps over semiring modules.\\
\addlinespace
Minimal connection API
& Lean-backed III
& \leanname{Transport.Connection}; \leanname{cocycleDefect}; \leanname{triangleHolonomy}; \leanname{triangleHolonomy_eq_refl_of_flat}
& Only self-normalization is assumed. Reverse inversion and the cocycle are not structure fields.\\
\addlinespace
Twisted three-point loop
& Lean-backed III
& \leanname{twistedThreePointConnection}; \leanname{twisted_triangleHolonomy_ne_refl}
& Concrete rational connection. The selected loop sends \(1\) to \(2\).\\
\addlinespace
Reversible additive defect
& Lean-backed III
& \leanname{ReversibleConnection}; \leanname{triangleHolonomy_sub_id}; \leanname{cocycleDefect_eq_zero_iff_triangleHolonomy_eq_refl}
& Heterogeneous module fibers, reverse-edge inversion, and an additive defect over a commutative ring.\\
\addlinespace
Reversible finite paths
& Lean-backed III
& \leanname{pathTransport_eq_direct_of_flat}; \leanname{indexedPathTransport_eq_direct_of_flat}; \leanname{indexedLoopTransport_eq_refl_of_flat}
& Finite pair-connection paths collapse to endpoints under algebraic flatness.\\
\addlinespace
Gauge covariance
& Lean-backed III
& \leanname{gauge_cocycleDefect}; \leanname{gauge_triangleHolonomy}; \leanname{gauge_flat_iff}
& Fiberwise algebraic conjugation, without smooth gauge theory.\\
\addlinespace
Flat versus canonical transport
& Lean-backed III
& \leanname{flatGauge}; \leanname{flat_isGaugeOfCanonical}
& A chosen base chart and ambient frames exhibit every flat pair connection as a gauge transform of canonical transport.\\
\addlinespace
Operation-tree locality
& Lean-backed III
& \leanname{OperationTree}; \leanname{OperationTree.eval}; \leanname{OperationTree.eval_congr_on_inputs}
& Arbitrary unary and binary semantics. Only dependence on named inputs is proved.\\
\addlinespace
Generic obstruction packet
& Lean-backed III
& \leanname{ObstructionPacket.not_isSolution}; \leanname{not_isSolution_iff_nonempty_obstructionPacket}
& Exact candidate-specific logical certificate, without global classification.\\
\addlinespace
Holonomy block obstruction
& Lean-backed III
& \leanname{obstruction_sum_blocks}; \leanname{obstructionBlocks_zero_iff}; \leanname{obstructionDetection_equivalences}
& Finite complete idempotent packet family on the base fiber. Orthogonality is not required by the compiled reconstruction.\\
\addlinespace
Packet witness and block gauge
& Lean-backed III
& \leanname{exists_packetWitness_of_obstructionBlock_ne_zero}; \leanname{obstructionBlock_gauge_of_commutes}; \leanname{obstructionBlock_gauge_eq_zero_iff}
& Witness extraction is classical. Blockwise covariance requires the gauge to commute with every packet projector.\\
\addlinespace
Heterogeneous matrix paths
& Lean-backed III
& \leanname{heterogeneousPacketPathExpansion}; \leanname{heterogeneousProduct_apply_eq_zero_of_no_supportedPath}
& Arbitrarily many different matrices over one common finite packet type. Fully typed changing-index chains compile only at lengths two and three.\\
\addlinespace
Arbitrary-function coordinates
& Lean-backed III
& \leanname{exactCoordinateMap_comp}; \leanname{exactCoordinateMap_naturality}
& Exact set-level equivalences and arbitrary functions, without a linearity or differentiability assumption.\\
\addlinespace
Colored and smooth trees
& Roadmap only
& roadmap \leanname{ColoredOperationTree.eval_eq_sum_packetLabellings}
& Dependent colors, multilinear packet expansion, derivatives, and Fa\`a di Bruno identities remain outside the contract.\\
\bottomrule
\end{longtable}
\normalsize

\subsection{Kernel-backed Chapter III surface}

The authoritative declaration-level contract is
\path{formalization/lean/LEAN_THEOREM_MAP_v1.0.1.md}.  It groups the compiled
surface as follows.
\begin{itemize}
\item \textbf{EdgeOperators and EdgeExtensions:} scalar packet matrices,
typed edge and word support, finite word/walk expansions, endpoint coefficient
sums, fixed-length reachability, orbit support, and common-index
heterogeneous products with declared supports.
\item \textbf{Transport and ReversibleConnections:} exact frame transport,
the minimal connection, heterogeneous reversible pair transport, additive
defect and triangle holonomy, finite paths, canonical connections, and the
fiberwise gauge action.
\item \textbf{CoordinateNaturality and HolonomyBlocks:} composition and
naturality for arbitrary functions between exact coordinate systems, finite
two-sided block reconstruction, equivalent obstruction criteria,
packet-supported witnesses, and packet-preserving block gauges.
\item \textbf{NonlinearObstructions:} finite unary--binary operation trees,
input locality, and candidate-specific obstruction certificates.
\end{itemize}
The formal companion records the checker-generated declaration
and audit totals.  The parity gate requires every fully qualified Chapter III
target printed beside a statement to occur in the theorem map, rejects stale
or duplicate targets, and checks it against the compiled environment.
\path{KernelAudit.lean} reports dependencies for the mapped and supporting
theorems, and its fail-closed gate scans every compiled
\leanname{Komorebi.PaperIII} declaration.  The four inherited statement
contracts remain identified separately by their inherited declaration names.

\begin{boundary}[No scope promotion by prose]
The displayed endpoint coefficient formula, reachability results, reversible
connection theorems, additive defect--holonomy identity, gauge covariance,
holonomy-block reconstruction, packet witnesses, common-index heterogeneous
path expansion, and arbitrary-function coordinate naturality are all part of
the finite algebraic kernel-backed surface.  What remains outside is the
dependent fiber-valued edge-system generalization, a fully dependent
arbitrary-length chain whose module and packet type change at every stage,
dependent colored operation trees, packet-labeled multilinear tree
expansions, nonlinear tree-defect and Fa\`a di Bruno identities, and the
analytic and geometric extensions listed above.
\end{boundary}

\section{The finite contract and the nonlinear transition}

Typed packet support turns the expansion of a transfer power into a walk
calculus.  Incompatible words vanish because their interfaces do not match.
Canonical packet coordinate transport is the exact equivalence obtained
from analysis and synthesis.  It satisfies the cocycle law, so every finite
canonical loop telescopes to the identity.  Holonomy enters after a separate
connection is supplied.  The compiled minimal API proves that flatness makes
both the relative defect and holonomy trivial, and it exhibits a nontrivial
rational loop.  Under the compiled reversible API, a nonzero additive triangle
defect is equivalent to nonidentity holonomy, transforms covariantly under
gauge, and is detected exactly by a finite family of packet blocks.  Finite
flat paths collapse to their endpoints.  Every flat reversible pair connection
is explicitly related to canonical transport by a gauge change through a
chosen base chart.

These facts define the boundary of the finite formal contract.  At that
boundary, exact coordinates preserve composition and satisfy naturality for
arbitrary functions in Chapter III, while the typed linear packet coordinate
counterpart remains inherited from Chapter~II.  Nontrivial nonlinear defects
must therefore be attached to additional modeling choices or to
packet compatibility of the particular operation.  Colored operation trees,
multilinear packet labelings, and Fa\`a di Bruno expansions remain
outside this finite Lean contract. Part~II develops the corresponding ordinary mathematics, including packet jets, rooted tree expansions, and explicit nonlinear obstruction theorems.


\clearpage
